% =====================================================================
%  2bat-ud2-12.tex  —  SESSIÓ 12: Casos especials (solució múltiple / regió no acotada)
% =====================================================================
\horatitol{Sessió 12 — Casos especials: solució múltiple o regió no acotada}%
{Dos casos que generen confusió: quan hi ha més d'una solució òptima i quan la regió s'estén indefinidament.}

\subsection{Idea clau}

El mètode (avaluar als vèrtexs) sempre funciona, però hi ha \textbf{dos casos
especials} que convé reconèixer per no confondre's. Avui els descobrim a les
pissarres.

\begin{idea}
\textbf{Cas 1 — solució múltiple.} Si la funció objectiu pren el \textbf{mateix valor
òptim} en \textbf{dos vèrtexs} veïns, aleshores \emph{tot el segment} que els uneix és
òptim: hi ha \textbf{infinites solucions}, totes igual de bones. Passa quan la funció
objectiu és «paral·lela» a una vora de la regió.

\textbf{Cas 2 — regió no acotada.} Si la regió factible s'estén indefinidament (no és
tancada), pot passar que:
\begin{itemize}[leftmargin=1.4em, itemsep=2pt]
  \item El \textbf{mínim} existeix i s'assoleix en un vèrtex (típic dels problemes de
        cost).
  \item El \textbf{màxim} \textbf{no existeix}: la funció pot créixer sense límit.
\end{itemize}
\end{idea}

\subsection{Exemples}

\begin{exemple}[solució múltiple]
Sobre la regió de vèrtexs $(0,0)$, $(6,0)$, $(3,3)$, $(0,4)$, maximitzem
$P=2x+2y$. Avaluem:
\[
P(0,0)=0,\quad P(6,0)=12,\quad P(3,3)=12,\quad P(0,4)=8.
\]
El màxim $P=12$ s'assoleix \textbf{a la vegada} a $(6,0)$ \emph{i} a $(3,3)$. De fet,
tot el segment que els uneix (la vora $x+y=6$) dóna $P=12$: hi ha \textbf{infinites}
solucions òptimes.

\begin{center}
\begin{tikzpicture}
\begin{axis}[udaxis, width=8cm, height=5cm,
    xlabel={\small $x$}, ylabel={\small $y$},
    xmin=0, xmax=7, ymin=0, ymax=5, xtick={0,2,4,6}, ytick={0,2,4}]
  \fill[udblue!16] (axis cs:0,0) -- (axis cs:6,0) -- (axis cs:3,3) -- (axis cs:0,4) -- cycle;
  \draw[udnavy, line width=1.6pt] (axis cs:6,0) -- (axis cs:3,3);
  \addplot[udnavy, only marks, mark=*, mark size=1.5pt] coordinates {(6,0)(3,3)};
  \node[font=\scriptsize, anchor=south west] at (axis cs:3,3) {$(3,3)$};
  \node[font=\scriptsize, anchor=north] at (axis cs:6,-0.1) {$(6,0)$};
  \node[udnavy, font=\scriptsize, anchor=south west] at (axis cs:4,1.5) {tot òptim};
\end{axis}
\end{tikzpicture}

\figcap{La vora ressaltada (de $(6,0)$ a $(3,3)$) és tota òptima: solució múltiple.}
\end{center}
\end{exemple}

\begin{exemple}[regió no acotada: hi ha mínim però no màxim]
En el cas de minimitzar cost de la sessió anterior, la regió era no acotada cap amunt.
El \textbf{mínim} de $C=40x+30y$ existia i era $C=170$ al vèrtex $(2,3)$. Però si
intentéssim \emph{maximitzar} $C$ sobre aquesta regió, \textbf{no hi hauria màxim}:
com que la regió s'estén indefinidament, podem agafar punts amb $x$ i $y$ tan grans
com vulguem, i $C$ creixeria sense límit. \textbf{Per això}, en regions no acotades,
té sentit buscar el mínim, però no sempre el màxim.
\end{exemple}

\subsection{Exercicis resolts}

\begin{exresolt}[reconèixer una solució múltiple]
Sobre la regió de vèrtexs $(0,0)$, $(4,0)$, $(2,2)$, $(0,3)$, maximitza $P=x+y$. Què
observes?
\solucio
Avaluem: $P(0,0)=0$, $P(4,0)=4$, $P(2,2)=4$, $P(0,3)=3$. El màxim $P=4$ s'assoleix
\textbf{a $(4,0)$ i a $(2,2)$}: és una \textbf{solució múltiple}. Tot el segment entre
aquests dos vèrtexs és òptim.
\resposta{Solució múltiple: $P=4$ a $(4,0)$ i $(2,2)$, i a tot el segment que els uneix.}
\end{exresolt}

\begin{exresolt}[interpretar socialment la solució múltiple]
En el cas anterior, què vol dir, en termes reals, que hi hagi infinites solucions
òptimes?
\solucio
Vol dir que hi ha \textbf{diverses maneres igualment bones} de repartir els recursos:
totes les combinacions del segment òptim donen el mateix valor de la funció objectiu
(per exemple, la mateixa cobertura o el mateix cost). L'entitat pot triar entre elles
segons \emph{altres} criteris (preferències, equitat, comoditat), perquè
matemàticament són equivalents.
\resposta{Hi ha diverses decisions igual d'òptimes; es pot triar entre elles per altres criteris.}
\end{exresolt}

\subsection{Exercicis per fer}

\begin{exercicis}
\begin{exlist}
  \item Sobre la regió de vèrtexs $(0,0)$, $(5,0)$, $(2,3)$, $(0,5)$, maximitza
        $P=2x+2y$. Hi ha una o més solucions òptimes? Quines?
  \item Sobre la mateixa regió, maximitza $P=3x+y$. Ara n'hi ha una de sola? On?
  \item Per a una regió no acotada amb vèrtexs de baix $(0,6)$, $(3,2)$, $(7,0)$,
        minimitza $C=2x+3y$. On és el mínim?
  \item Sobre la regió de l'exercici 3, té sentit buscar el \textbf{màxim} de
        $C=2x+3y$? Justifica-ho.
  \item \textbf{(Reconèixer el cas)} Donada una regió i una funció objectiu, com pots
        \emph{sospitar}, abans de calcular, que hi haurà solució múltiple? (Pista:
        pensa en el pendent de la funció objectiu i el de les vores.)
  \item \textbf{(Síntesi)} Completa la frase amb les teves paraules: «En una regió no
        acotada, el \rule{1.5cm}{0.4pt} sempre existeix i és en un vèrtex, però el
        \rule{1.5cm}{0.4pt} pot no existir.»
\end{exlist}
\end{exercicis}

% ---------------------------------------------------------------------
\fullsolucions{Sessió 12}
\begin{solucions}
\begin{sollist}
  \item $P=2x+2y$: $P(0,0)=0$, $P(5,0)=10$, $P(2,3)=10$, $P(0,5)=10$. El màxim $P=10$
        s'assoleix a \textbf{tres} vèrtexs ($(5,0)$, $(2,3)$, $(0,5)$): solució múltiple
        (de fet, tota la vora superior).
  \item $P=3x+y$: $P(0,0)=0$, $P(5,0)=15$, $P(2,3)=9$, $P(0,5)=5$. Màxim únic $P=15$
        a $(5,0)$.
  \item $C(0,6)=18$, $C(3,2)=6+6=12$, $C(7,0)=14$. Mínim $C=12$ a $(3,2)$.
  \item No: la regió és no acotada cap amunt, així que $C=2x+3y$ pot créixer sense
        límit; el \textbf{màxim no existeix}.
  \item Resposta oberta. Idea: si la funció objectiu té el \emph{mateix pendent} que
        una de les vores de la regió, serà paral·lela a aquella vora i tot el segment
        donarà el mateix valor (solució múltiple).
  \item «El \textbf{mínim} sempre existeix i és en un vèrtex, però el \textbf{màxim}
        pot no existir.»
\end{sollist}
\end{solucions}
